\section{Diophantine equations, continued fractions, and bounded approximation}
\label{sec:diophantine}

This chapter uses signed integers, positive denominators, finite lists, and
bounded search only. A rational is represented by $(a,b)$ with $a\in\Z$
and $b\geq1$; equality and order are cross-multiplication tests.
Rational inequalities therefore abbreviate integer inequalities.
An algorithmic square root means the integer function
\[
 \operatorname{isqrt}(n)=\max\{j\leq n:j^2\leq n\}.
\]
Its specification $s^2\leq n<(s+1)^2$ follows by a bounded scan,
including $n=0$. No real-number decision procedure is assumed.

\subsection{Rational arithmetic and power obstructions}

\begin{theorem}[Canonical rational representation]\label{dio:normal}
Every pair $(a,b)$ with $b>0$ has a unique equivalent pair $(u,v)$ with
$v>0$ and $\gcd(|u|,v)=1$. Normalization, addition, multiplication,
comparison, and reciprocal of a nonzero rational are primitive recursive.
\end{theorem}
\begin{proof}
Put $g=\gcd(|a|,b)$, $u=a/g$, and $v=b/g$. Since $b>0$, $g>0$;
the quotients are exact. A common divisor larger than $1$, multiplied
by $g$, would contradict the greatest-common-divisor property.
If $u/v=u'/v'$ and both pairs are reduced, then $uv'=u'v$.
The coprime version of Euclid's lemma gives $v\mid v'$ and $v'\mid v$.
Positivity gives $v=v'$, and cancellation gives $u=u'$.
Addition and multiplication form
$(uv'+u'v,vv')$ and $(uu',vv')$ and normalize.
Comparison tests $uv'\leq u'v$. The reciprocal switches numerator
and denominator and moves a negative sign to the numerator.
Every operation uses integer arithmetic and gcd. Zero becomes $(0,1)$.
\end{proof}

\begin{theorem}[Rational root test]\label{dio:rational-root}
Let $f(X)=c_0+c_1X+\cdots+c_dX^d$, with $d\geq1$, $c_d\ne0$,
and integer coefficients. If a reduced rational $u/v$, $v>0$, is a
root, then $v\mid c_d$; if $u\ne0$, then $u\mid c_0$.
A rational root of a monic integer polynomial is therefore an integer.
If $c_0\ne0$, all rational roots can be found by a finite search with
$|u|\leq|c_0|$ and $1\leq v\leq|c_d|$.
\end{theorem}
\begin{proof}
Multiply the equation by $v^d$:
\[
 c_0v^d+c_1uv^{d-1}+\cdots+c_du^d=0.
\]
Reduction modulo $v$ shows $v\mid c_du^d$. Coprimality gives
$\gcd(u^d,v)=1$, hence $v\mid c_d$. Reduction modulo $u$ similarly
gives $u\mid c_0v^d$, hence $u\mid c_0$. For $c_d=1$ we obtain $v=1$.
If $c_0\ne0$, zero is not a root, and the divisibilities bound every
candidate. Horner evaluation after clearing denominators checks each
candidate by an integer equality. The degree and coefficient list are
finite data, so the complete test is primitive recursive.
\end{proof}

\begin{corollary}[Rational powers and irrational roots]\label{dio:powers}
For $N\in\N$ and $k\geq2$, a nonnegative rational whose $k$th power is
$N$ is an integer. Such a rational exists precisely when
$N=t^k$ for some $0\leq t\leq N+1$.
For $N>0$, the equivalent valuation criterion is
$k\mid v_p(N)$ for every prime $p\leq N$.
In particular, if $D$ is nonsquare and $q>0$, then
$p^2-Dq^2\ne0$ for every integer $p$.
\end{corollary}
\begin{proof}
Apply the rational root test to $X^k-N$. An integer root is bounded
by $N+1$, including $N=0$. For $N>0$, unique factorization makes
$t^k=N$ equivalent to every prime exponent being divisible by $k$;
the reverse implication uses the finite product
$\prod_{p\leq N}p^{v_p(N)/k}$.
If $p^2=Dq^2$, normalize $|p|/q$ and apply the square case.
Thus the irrationality of $\sqrt D$ has the purely arithmetical content
$q>0\Longrightarrow p^2\ne Dq^2$.
\end{proof}

\subsection{Pythagorean triples and two finite descents}

\begin{theorem}[Primitive Pythagorean triples]\label{dio:pythagoras}
Suppose $a,b,c>0$, $a^2+b^2=c^2$, and $\gcd(a,b)=1$.
Exactly one leg is even. Calling it $b$, there are unique integers
$m>n>0$ such that
\[
 a=m^2-n^2,\quad b=2mn,\quad c=m^2+n^2,\quad
 \gcd(m,n)=1,\quad m\not\equiv n\pmod2.
\]
Conversely these conditions produce a primitive positive triple.
The parameters can be recovered by bounded square-root tests.
\end{theorem}
\begin{proof}
The legs cannot both be even by coprimality. They cannot both be odd,
because their squared sum would be $2$ modulo $4$. Thus $b$ is even
and $a,c$ odd. A prime dividing $a,c$ divides $b^2$, hence $b$;
therefore $\gcd(a,c)=1$. The positive integers
\[
 U=(c+a)/2,\qquad V=(c-a)/2
\]
satisfy $U>V$, $UV=(b/2)^2$, and $\gcd(U,V)=1$: a common divisor
divides $c$ and $a$. In a coprime product that is a square, each prime
exponent in each factor is even. Hence $U=m^2$ and $V=n^2$ for unique
positive $m,n$. The displayed formulas follow. Coprimality follows
from that of $U,V$, and opposite parity follows because $c$ is odd.

Conversely expand
$(m^2-n^2)^2+(2mn)^2=(m^2+n^2)^2$.
The first leg is odd. An odd prime dividing both legs divides $m$ or
$n$, and then the other, contradicting coprimality.
The positive square roots of $U,V$ prove uniqueness; they are bounded
by $c$, so recovery is primitive recursive.
\end{proof}

\begin{corollary}[All integer right triangles]\label{dio:all-triples}
Every positive integer right triangle is a positive integer multiple of
exactly one primitive triple, up to exchanging its legs.
\end{corollary}
\begin{proof}
Put $g=\gcd(a,b)$. Since $g^2\mid c^2$, comparison of prime
exponents gives $g\mid c$. Divide by $g$ and apply the preceding
theorem. Conversely scaling preserves the equation.
The scale is forced to equal the gcd of the two legs.
\end{proof}

\begin{theorem}[No integer right triangle has positive square area]
\label{dio:square-area}
There are no positive integers $a,b,c,t$ satisfying
$a^2+b^2=c^2$ and $ab=2t^2$.
\end{theorem}
\begin{proof}
We give a size-decreasing transformation with explicitly bounded
iteration. First divide a proposed triangle by $g=\gcd(a,b)$.
The primitive triangle has integral area $A_0$, since one leg is even.
Then $t^2=g^2A_0$, so $g\mid t$ by prime exponents; its area is the
square $(t/g)^2$. This normalization does not increase the hypotenuse.

For a primitive counterexample the area, in the parameters of
Theorem~\ref{dio:pythagoras}, is
\[
 mn(m+n)(m-n).
\]
The four factors are pairwise coprime. Coprimality involving $m$ or
$n$ follows from $\gcd(m,n)=1$; a common divisor of $m+n,m-n$
divides $2m,2n$, and these two factors are odd.
Thus each factor is a square, say
\[
 m=u^2,\quad n=v^2,\quad m+n=w^2,\quad m-n=z^2.
\]
Both $w,z$ are odd. Since $w^2-z^2=2v^2$ and odd squares are $1$
modulo $8$, $v$ is even. Put $r=(w+z)/2$, $s=(w-z)/2$.
These are positive integers, and
\[
 r^2+s^2=(w^2+z^2)/2=m=u^2,\qquad
 rs/2=(w^2-z^2)/8=(v/2)^2.
\]
So $(r,s,u)$ is another integer right triangle of square area, with
smaller hypotenuse because $u^2=m<m^2+n^2=c$.

The step uses gcd, exact division, and bounded square-root searches.
Make it a total primitive recursive function by returning a fixed
default tuple on invalid input. Starting at hypotenuse $c_0$, iterate
it $c_0$ times. The quantifier-free invariant is that the current tuple
is a counterexample with $1\leq c_i\leq c_0-i$. At $i=c_0$ this is
impossible. This uses quantifier-free induction, not an unrestricted
least-counterexample principle.
\end{proof}

\begin{theorem}[Fermat's quartic descent]\label{dio:fermat-four}
There are no positive integers $x,y,z$ with $x^4+y^4=z^2$.
Consequently $x^{4k}+y^{4k}=z^{4k}$ has no positive integer solution
for any $k\geq1$.
\end{theorem}
\begin{proof}
Normalize by $g=\gcd(x,y)$. The equation gives $g^4\mid z^2$,
hence $g^2\mid z$, so division yields a solution with coprime $x,y$
and no larger $z$. They cannot both be odd, by reduction modulo $16$.
Exchange them so $x$ is odd and $y$ even.
The primitive right triangle $(x^2,y^2,z)$ has parameters
\[
 x^2=m^2-n^2,\qquad y^2=2mn,\qquad z=m^2+n^2.
\]
If $m$ were even and $n$ odd, the first equality would make a square
$3$ modulo $4$. Hence $m$ is odd and $n$ even.
The coprime factors $m,2n$ have square product, so $m=u^2$ and
$2n=w^2$. The latter implies $w=2v$ and $n=2v^2$.

Now $(x,n,m)$ is a primitive right triangle: a prime dividing $x,n$
would also divide $m$. Its even leg is $2v^2$ and its hypotenuse is
$u^2$. Another parametrization gives coprime positive $r>s$ with
\[
 x=r^2-s^2,\qquad v^2=rs,\qquad u^2=r^2+s^2.
\]
Therefore $r=A^2,s=B^2$, and $A^4+B^4=u^2$ is a new positive
solution. It is smaller since $u^2=m<z=m^2+n^2$.
Normalize and apply this transformation with fuel equal to the original
$z_0$. The quantifier-free invariant records the quartic equation and
$1\leq z_i\leq z_0-i$, giving a contradiction at the last stage.
All root searches are bounded by the current hypotenuse.
For the last assertion substitute $X=x^k,Y=y^k,Z=z^{2k}$.
\end{proof}

\subsection{Finite continued fractions and Farey arithmetic}

For $a_0\geq0$ and $a_i\geq1$ for $i>0$, write
$[a_0;a_1,\ldots,a_n]$ for the usual right-nested rational fraction.
Define continuants by
\[
 \begin{gathered}
 p_{-2}=0,\quad p_{-1}=1,\quad
 q_{-2}=1,\quad q_{-1}=0,\\
 p_i=a_ip_{i-1}+p_{i-2},\qquad
 q_i=a_iq_{i-1}+q_{i-2}\quad(i\geq0).
 \end{gathered}
\]
Negative subscripts name four initial constants; an implementation
stores a four-entry state.

\begin{theorem}[Finite continued fractions for rationals]\label{dio:cf-rational}
Every nonnegative rational has a unique finite simple continued fraction
whose last coefficient is at least $2$, unless there is just one
coefficient. For $a/b$, $a\geq0,b>0$, it can be computed in at most
$b+1$ Euclidean stages, with every coefficient at most $a+b$.
\end{theorem}
\begin{proof}
Divide $a=a_0b+r$. If $r=0$, stop; otherwise replace $(a,b)$ by
$(b,r)$ and repeat. Every nonterminal positive second entry decreases,
so $b+1$ stages suffice. Substitution backwards gives the continued
fraction. Later quotients are positive. At a final stage after the
first, the first entry is larger than, and a multiple of, the second,
so the quotient is at least $2$. All entries are at most $a+b$.
For uniqueness, the tail of a nonintegral canonical fraction is greater
than $1$, so its first coefficient is the floor of its value.
Subtract the coefficient, invert the positive remainder, and repeat.
This recovers every coefficient. A final $1$ accounts exactly for the
excluded alternate expression $[\ldots,c,1]=[\ldots,c+1]$.
All floors here are rational floors, computed by integer division.
\end{proof}

\begin{theorem}[Continuant identity, coprimality, and growth]
\label{dio:continuants}
For every valid finite coefficient list,
\[
 [a_0;\ldots,a_n]=p_n/q_n,\qquad
 p_nq_{n-1}-p_{n-1}q_n=(-1)^{n+1}.
\]
In particular $\gcd(p_n,q_n)=1$ and $q_n>0$.
The denominators are nondecreasing from $q_0=1$ and strictly increase
after $q_1$. Whenever the indicated terms exist,
\[
 q_{n+2}\geq2q_n,\qquad q_{2k}\geq2^k.
\]
If every coefficient is at most $A\geq1$, then
$p_n,q_n\leq(A+1)^{n+1}$.
\end{theorem}
\begin{proof}
Induction by matrix multiplication gives
\[
 \begin{pmatrix}a_0&1\\1&0\end{pmatrix}\cdots
 \begin{pmatrix}a_n&1\\1&0\end{pmatrix}
 =
 \begin{pmatrix}p_n&p_{n-1}\\q_n&q_{n-1}\end{pmatrix}.
\]
Substitution in the nested fraction gives its value $p_n/q_n$.
Each factor has determinant $-1$, proving the determinant identity.
Any common divisor of $p_n,q_n$ divides $1$.
The recurrence gives $q_0=1,q_1=a_1\geq1$, and
$q_{n+2}\geq q_{n+1}+q_n>q_{n+1}$.
Thus $q_{n+2}\geq2q_n$, and induction gives the even-index bound.
For the upper bound start with $p_0\leq A,q_0=1$ and
$p_1\leq A^2+1\leq(A+1)^2,q_1\leq A$. Thereafter use
\[
 A(A+1)^n+(A+1)^{n-1}\leq(A+1)^{n+1}.
\]
All these are invariants of a finite matrix-product loop.
\end{proof}

\begin{theorem}[Exact tail and error identities]\label{dio:cf-error}
For a positive rational $t$ and $n\geq0$,
\[
 [a_0;\ldots,a_n,t]
 =\frac{tp_n+p_{n-1}}{tq_n+q_{n-1}},\qquad
 \left|[a_0;\ldots,a_n,t]-\frac{p_n}{q_n}\right|
 =\frac1{q_n(tq_n+q_{n-1})}.
\]
For $t\geq1$ the distance from $p_n/q_n$ is at most $1/q_n^2$.
For $n\geq1$ the value lies between adjacent convergents.
If a next coefficient $a_{n+1}$ is specified and $t\geq a_{n+1}$,
its distance is at most $1/(q_nq_{n+1})$.
\end{theorem}
\begin{proof}
Insert $t$ in the matrix identity and subtract $p_n/q_n$.
The numerator of the difference is
$p_{n-1}q_n-p_nq_{n-1}$, of absolute value $1$.
All denominators are positive. For $n\geq1$ the value is the weighted
mean of $p_n/q_n,p_{n-1}/q_{n-1}$ with positive weights
$tq_n,q_{n-1}$. Finally $tq_n+q_{n-1}\geq q_n$ if $t\geq1$,
and is at least $q_{n+1}$ under the stronger hypothesis.
For a finite extension $t$ is a positive rational, so all claims are
finite integer identities after clearing denominators.
\end{proof}

\begin{corollary}[A primitive recursive continued-fraction real]
\label{dio:cf-real}
Fix a primitive recursive coefficient function $a(n)$ with $a(0)\geq0$
and $a(n)\geq1$ for $n>0$. Its convergents satisfy
\[
 m>n\quad\Longrightarrow\quad
 |p_m/q_m-p_n/q_n|\leq q_n^{-2}.
\]
Consequently $r_k=p_{2k+2}/q_{2k+2}$ is a primitive recursive rational
name with $|r_\ell-r_k|\leq2^{-2k-2}$ whenever $\ell>k$.
\end{corollary}
\begin{proof}
For specified finite $m>n$, use the preceding theorem with
$t=[a_{n+1};\ldots,a_m]\geq1$. Then
$q_{2k+2}\geq2^{k+1}$ gives the stated estimate.
Each prefix is produced by primitive recursion.
This is a schema for each fixed primitive recursive coefficient
function and its positivity proofs; it does not say that arbitrary
reals have primitive recursive continued-fraction expansions.
The estimates are themselves the arithmetical Cauchy data, without
completion or choice.
\end{proof}

\begin{theorem}[Farey separation and the mediant]\label{dio:farey}
Suppose $a/b<c/d$, $b,d>0$, and $bc-ad=1$.
Every rational $p/q$ strictly between them, $q>0$, has $q\geq b+d$.
The mediant $(a+c)/(b+d)$ is reduced, lies strictly between them,
and attains that denominator bound.
\end{theorem}
\begin{proof}
The strict inequalities give $bp-aq\geq1$ and $cq-dp\geq1$.
Hence
\[
 q=d(bp-aq)+b(cq-dp)\geq d+b.
\]
Cross multiplication places the mediant strictly between the endpoints.
A common divisor of its numerator and denominator divides
$b(a+c)-a(b+d)=bc-ad=1$.
In particular this applies to appropriately ordered adjacent convergents.
\end{proof}

\subsection{Finite Dirichlet approximation}

\paragraph{The bounded collision principle.}
We use Theorem~\ref{fd:php}: a list of $M+1$ integers in
$\{0,\ldots,M-1\}$ contains an equal pair, which can be found by
checking at most $(M+1)^2$ pairs.
For the present applications, one can maintain a Boolean array of
length $M$ marking encountered values.
Processing a new distinct value changes exactly one zero to one.
After $j$ distinct entries the array has exactly $j$ ones, by induction.
It cannot contain $M+1$ ones. Thus the scan finds a collision.
Array counts and tests for repeated entries in a finite prefix are
bounded Boolean operations; their induction invariants are quantifier-free.

\begin{theorem}[Rational Dirichlet approximation]\label{dio:dirichlet}
For a rational $\alpha$ and $Q\geq1$, there are integers $p,q$ with
\[
 1\leq q\leq Q,\qquad |q\alpha-p|<1/Q.
\]
They can be found by a bounded scan of $Q+1$ rational fractional parts.
\end{theorem}
\begin{proof}
For $j=0,\ldots,Q$, let
$f_j=j\alpha-\lfloor j\alpha\rfloor\in[0,1)$ and assign it box
$\lfloor Qf_j\rfloor<Q$.
Two indices $i<j$ have equal box numbers. Put
$q=j-i$ and $p=\lfloor j\alpha\rfloor-\lfloor i\alpha\rfloor$.
Then $q\alpha-p=f_j-f_i$, of absolute value less than $1/Q$.
For $\alpha=a/b$, $b>0$, compute the fractional parts from
remainders of $ja$ modulo $b$. There is no real floor.
The witness satisfies $|p|\leq Q|\alpha|+1$.
\end{proof}

\begin{theorem}[Simultaneous rational Dirichlet approximation]
\label{dio:simultaneous}
For a finite rational list $\alpha_1,\ldots,\alpha_d$, $d\geq1$,
and $N\geq1$, there are $q,p_1,\ldots,p_d$ satisfying
\[
 1\leq q\leq N^d,\qquad |q\alpha_i-p_i|<1/N
 \quad(1\leq i\leq d).
\]
The construction and witness bounds are primitive recursive in the
finite input list and $N$.
\end{theorem}
\begin{proof}
For $j=0,\ldots,N^d$, form the vector of box numbers
$\lfloor N\{j\alpha_i\}\rfloor<N$. Encode it in base $N$;
for $N=1$ every vector has the single code zero.
There are $N^d$ codes and $N^d+1$ vectors. A collision at $j>k$,
with $q=j-k$, gives the result by subtracting the integer parts in
every coordinate. All errors are less than the box width $1/N$.
One may use $|p_i|\leq N^d|\alpha_i|+1$.
Exponentiation, finite-list operations, and the collision scan are
primitive recursive.
\end{proof}

\begin{theorem}[Certified transfer to represented reals]\label{dio:real-transfer}
Fix rational approximation names for a finite list of quantities
$\xi_1,\ldots,\xi_d$. Suppose a chosen precision supplies rationals
$r_i$ with certificates $|\xi_i-r_i|\leq\varepsilon$.
The rational construction produces $1\leq q\leq N^d$ and integers
$p_i$ with
\[
 |q\xi_i-p_i|<1/N+N^d\varepsilon.
\]
In particular $\varepsilon\leq1/(2N^{d+1})$ gives errors
less than $3/(2N)$.
\end{theorem}
\begin{proof}
Take witnesses for the rational list $r_i$. Then
\[
 |q\xi_i-p_i|\leq q|\xi_i-r_i|+|qr_i-p_i|
 <N^d\varepsilon+1/N.
\]
Arithmetically, replace $\xi_i$ by all later rational approximants and
the supplied error inequalities; the same calculation proves each
finite inequality, uniformly in the later precision index.
For names with error $2^{-k}$, the generous choice $k=2N^{d+1}$
suffices, since $2^k\geq k$ by quantifier-free induction.
Thus the required precision is primitive recursive.
This is a schema for fixed names and their approximation proofs.
It retains explicit slack; it neither assumes a decidable real floor
nor claims the sharp real version with the identical denominator bound.
\end{proof}

\subsection{Pell's equation with an explicit primitive recursive bound}

The following proof uses rational approximation and a finite collision.
It needs neither periodic continued fractions nor a limiting argument.

\begin{lemma}[Bounded norm with a prescribed lower denominator]
\label{dio:pell-approx}
Let $D\geq2$ be nonsquare and $Y\geq0$. Set
\[
 B=4(D+1),\qquad Q=B(Y+1),\qquad C=4D+2.
\]
There are positive integers $p,q$ with
\[
 Y<q\leq Q,\qquad p\leq DQ+1,\qquad 0<|p^2-Dq^2|<C.
\]
The pair can be obtained by bounded search using integer arithmetic.
\end{lemma}
\begin{proof}
Put $A=\operatorname{isqrt}(DQ^4)$ and $\alpha=A/Q^2$. Then
\[
 1\leq\alpha\leq D,\qquad
 0\leq D-\alpha^2<\frac{2A+1}{Q^4}
 \leq\frac{2D}{Q^2}+\frac1{Q^4}.
\]
Indeed $A\geq Q^2$, since $D\geq2$, and
$A^2\leq DQ^4\leq D^2Q^4$.
Rational Dirichlet approximation gives $1\leq q\leq Q$ and $p$ with
$|p-q\alpha|<1/Q$. Since $Q\geq12$ and $\alpha\geq1$, $p>0$.
Also $p\leq Dq+1\leq DQ+1$. Factoring a difference of squares,
\[
 \begin{split}
 |p^2-Dq^2|
 &\leq |p-q\alpha|(p+q\alpha)+q^2(D-\alpha^2)\\
 &<\frac{2Dq+1}{Q}
       +q^2\left(\frac{2D}{Q^2}+\frac1{Q^4}\right)\\
 &\leq4D+\frac1Q+\frac1{Q^2}<C.
 \end{split}
\]
It is nonzero by Corollary~\ref{dio:powers}.
If $q\leq Y$, the same estimate gives
\[
 |p^2-Dq^2|
 <\frac{2DY+1}{4(D+1)(Y+1)}
  +\frac{(2D+1)Y^2}{16(D+1)^2(Y+1)^2}<1.
\]
The first summand is less than $1/2$, since
$2DY+1\leq2D(Y+1)$; the second is less than $1/2$, since
$2D+1<8(D+1)^2$. This contradicts the nonzero integral norm.
Thus $q>Y$. Clearing the positive denominators expresses the entire
proof as integer equalities and inequalities.
\end{proof}

\begin{theorem}[Uniform bounded Pell existence]\label{dio:pell-existence}
For nonsquare $D\geq2$, define
\[
 \begin{gathered}
 C=4D+2,\quad B=4(D+1),\quad K=2C^3,\\
 U=(B+1)^{K+1},\qquad H=4(D+1)^2(U+1)^2.
 \end{gathered}
\]
There are integers
\[
 2\leq x\leq H,\qquad1\leq y\leq H,\qquad x^2-Dy^2=1.
\]
The first such pair in the displayed finite box is therefore a primitive
recursive function of $D$ on nonsquare inputs.
\end{theorem}
\begin{proof}
Let $T_0=0,T_{i+1}=B(T_i+1)$. For each $i=0,\ldots,K$ use the
preceding lemma with $Y=T_i$ to obtain $(p_i,q_i)$ and
$m_i=p_i^2-Dq_i^2$, with
\[
 T_i<q_i\leq T_{i+1},\qquad 0<|m_i|\leq C,\qquad
 p_i\leq DT_{i+1}+1.
\]
The $q_i$ strictly increase. Induction gives
$T_i+1\leq(B+1)^i$, so $q_i<U$ and $p_i\leq DU+1$.
Classify a pair by
\[
 (m_i,\ p_i\bmod|m_i|,\ q_i\bmod|m_i|).
\]
There are at most $2C^3=K$ classes: use a sign-and-magnitude tag
less than $2C$, and two residue slots each less than $C$.
A collision among the $K+1$ pairs gives $i<j$ with common nonzero
norm $m$ and matching residues modulo $|m|$.
Define
\[
 x=\frac{p_ip_j-Dq_iq_j}{m},\qquad
 y=\frac{p_iq_j-p_jq_i}{m}.
\]
The first numerator is congruent to
$p_i^2-Dq_i^2=m$ modulo $|m|$, and the second to zero.
Thus both quotients are integers, also for $m<0$.
Expansion gives
\[
 (p_ip_j-Dq_iq_j)^2-D(p_iq_j-p_jq_i)^2=m^2,
\]
so their norm is $1$.

Positivity needs checking, especially for negative common norm.
Since $q_j>q_i>0$ and all $p_i$ are positive,
\[
 (p_iq_j)^2-(p_jq_i)^2=m(q_j^2-q_i^2).
\]
The second numerator has the same sign as $m$, hence $y>0$.
Also
\[
 (p_ip_j)^2-(Dq_iq_j)^2
 =m\bigl(D(q_i^2+q_j^2)+m\bigr).
\]
The bracket is positive: for $m<0$ use
$Dq_i^2+m=p_i^2>0$, and for $m>0$ it is immediate.
The first numerator therefore has the same sign as $m$, so $x>0$.
Now $D\geq2,y>0$, and the norm equation imply $x\geq2$.
Finally, because $|m|\geq1$,
\[
 x\leq(DU+1)^2+DU^2\leq H,\qquad
 y\leq2(DU+1)U\leq H.
\]
The deliberately loose estimates are sufficient.
There are exactly $K+1$ stages and each inner construction is bounded.
Alternatively the proved existence permits a direct scan of the
$H^2$ candidate pairs. Give the function a default value when
$D<2$ or $\operatorname{isqrt}(D)^2=D$ to make it total.
\end{proof}

\begin{theorem}[Pell composition and arbitrarily many solutions]
\label{dio:pell-powers}
Define
\[
 (x,y)\mathbin{\star}(u,v)=(xu+Dyv,xv+yu),\qquad
 \mathcal N_D(x,y)=x^2-Dy^2.
\]
The norm is multiplicative. If the norm is $1$, $(x,-y)$ is an inverse.
For $u>1,v>0,u^2-Dv^2=1$, define
\[
 (X_0,Y_0)=(1,0),\qquad
 X_{n+1}=uX_n+DvY_n,\quad Y_{n+1}=vX_n+uY_n.
\]
Every pair has norm $1$ and $Y_{n+1}>Y_n$. For every specified $n$,
this primitive recursive procedure gives $n$ distinct positive solutions.
\end{theorem}
\begin{proof}
Polynomial expansion proves norm multiplicativity, associativity,
and the conjugate inverse identity. Induction proves norm $1$ and
$X_n\geq1,Y_n\geq0$. The recurrence gives
$Y_{n+1}=vX_n+uY_n>Y_n$, since $vX_n\geq1,u\geq2$.
For a size bound put $L=u+Dv+1$. Induction gives
$X_n,Y_n\leq L^n$. The infinitude assertion is this finite witness
construction for each $n$, not an assertion that an infinite set exists.
\end{proof}

\begin{theorem}[All positive Pell solutions are fundamental powers]
\label{dio:pell-all}
Let $(u,v)$ be a norm-one solution with $u>1,v>0$ whose $v$ is
smallest among positive solutions. Such a pair is found in the box of
Theorem~\ref{dio:pell-existence}. Every solution with $x>0,y\geq0$
is uniquely $(u,v)^{\star n}$ for an integer $0\leq n\leq y$.
Its exponent is recovered by a primitive recursive procedure.
\end{theorem}
\begin{proof}
For $y=0$ the equation gives $x=1$. For $y>0$, minimality gives
$y\geq v$. Divide by $(u,v)$ using the integer pair
\[
 x'=ux-Dvy,\qquad y'=uy-vx.
\]
Its norm is $1$. Furthermore
\[
 (uy)^2-(vx)^2=y^2-v^2\geq0,
\]
so $y'\geq0$. Positivity of $x'$ follows from
\[
 (ux)^2-(Dvy)^2=1+D(v^2+y^2)>0.
\]
Finally
\[
 (vx)^2-\bigl((u-1)y\bigr)^2=v^2+2(u-1)y^2>0,
\]
so $vx>(u-1)y$, which says $y'<y$.
Repeat until reaching $(1,0)$, using at most the original $y$ steps.
The invariant states that the current pair has norm $1$, has
nonnegative coordinates with positive first coordinate, and multiplied
by the already recorded fundamental power equals the original pair.
The second coordinate supplies the decreasing fuel.
Reconstruction gives the power representation.
The preceding theorem makes second coordinates of distinct powers
strictly increasing, proving uniqueness.

For the fundamental solution, first find $(x_0,y_0)$ in the box of size
$H$. Scan $1\leq y\leq y_0,2\leq x\leq x_0$.
Any competitor with smaller $y$ has smaller $x$ by $x^2=1+Dy^2$,
so this scan finds a globally minimal positive $y$.
\end{proof}

\begin{corollary}[An explicit Pell family]\label{dio:pell-family}
For $t\geq2$, the nonsquare $D=t^2-1$ has fundamental solution
$(t,1)$. All its positive solutions are powers of this pair.
\end{corollary}
\begin{proof}
The inequalities $(t-1)^2<t^2-1<t^2$ prove nonsquareness.
The identity $t^2-(t^2-1)\cdot1^2=1$ gives a solution with smallest
possible positive second coordinate. Apply the preceding theorem.
\end{proof}

\subsection{Induction audit and Lean implementation route}

Each construction has a total primitive recursive version using a bounded
scan or fixed-fuel recurrence and a default on malformed inputs.
Rational claims are proved after clearing positive denominators.
Bounded universal conditions, such as preservation of previous entries,
are Boolean folds and can be part of a quantifier-free invariant.
The descents use the original positive size as fuel. Pell existence uses
$K+1$ stages and the explicit bound $H$; Pell classification uses the
input second coordinate.

Formalization has three distinct products. First define the algorithms
and prove their ordinary Lean specifications over naturals, integers, and
explicit rational pairs. Second give primitive recursive codes and
defining equations in the chosen iPRA presentation. Third construct
iPRA derivations of the invariants and specifications.
An ordinary Lean arithmetic proof alone does not complete the third step.

Useful module boundaries are \texttt{RationalCodes},
\texttt{PythagoreanParameters}, \texttt{BoundedDescent},
\texttt{Continuants}, \texttt{FiniteDirichlet},
\texttt{PellNormCollision}, and \texttt{PellClassification}.
The separate blueprint provides concrete definitions and proof obligations.
ProveIt modules named \texttt{Pythagoras}, \texttt{FermatQuartic},
\texttt{SquareArea}, \texttt{Dirichlet}, \texttt{SimultaneousGrid}, and
\texttt{PellGeneral} are integration points subject to the repository
audit; their names alone do not certify the explicit uniform bound above.
